\documentclass[a4paper, noanswers, addpoints, 12pt]{exam}
\usepackage[T1]{fontenc}
\usepackage[utf8]{inputenc}
\usepackage{lmodern}
\usepackage[french]{babel}
\usepackage{xcolor}
\usepackage{listings}
\usepackage{graphicx}
\usepackage{array}
\usepackage{eurosym}

% Configuration de la coloration syntaxique
\definecolor{codegreen}{rgb}{0,0.6,0}
\definecolor{codegray}{rgb}{0.5,0.5,0.5}
\definecolor{codepurple}{rgb}{0.58,0,0.82}
\definecolor{backcolour}{rgb}{0.95,0.95,0.92}
 
\lstdefinestyle{mystyle}{
    backgroundcolor=\color{backcolour},   
    commentstyle=\color{codegreen},
    keywordstyle=\color{magenta},
    numberstyle=\tiny\color{codegray},
    stringstyle=\color{codepurple},
    basicstyle=\footnotesize\ttfamily,
    breakatwhitespace=false,         
    breaklines=true,                 
    captionpos=b,                    
    keepspaces=true,                 
    numbers=left,                    
    numbersep=5pt,                  
    showspaces=false,                
    showstringspaces=false,
    showtabs=false,                  
    tabsize=2
}
\lstset{style=mystyle, language=Python}

\title{\vspace{-2cm}Examen Première NSI (Rattrapage) - Types construits et gloutonnerie}
\date{}

\usepackage[a4paper,top=2.3cm,bottom=2.8cm,left=2cm,right=2cm,headsep=0.6cm,footskip=1.2cm]{geometry}
\lstset{extendedchars=true,literate={û}{{\^u}}1}
\begin{document}
\runningfooter{}{\footnotesize Babash - NSI Tani Malandi et Johanne Bakalara \textperiodcentered{} CC BY-SA 4.0 \textperiodcentered{} page \thepage}{}
\firstpagefooter{}{\footnotesize Babash - NSI Tani Malandi et Johanne Bakalara \textperiodcentered{} CC BY-SA 4.0 \textperiodcentered{} page \thepage}{}

\maketitle
\thispagestyle{headandfoot}

\pagestyle{headandfoot}
\thispagestyle{headandfoot}
\firstpageheader{\makebox[\textwidth]{Nom : \dotfill\ Pr\'enom : \dotfill\ Classe : \makebox[3cm]{\dotfill}}}{}{}
\vspace{-3cm}

\section*{Consignes:}
\begin{center} 
\fbox{\parbox{\linewidth}{
\begin{center}
La clarté de la rédaction et les explications font partie intégrante de la notation. 
Pour les questions demandant du code, respectez la syntaxe Python. 
\end{center}
}}
\end{center}

\section*{Barème :}
\begin{center}
  \par\bigskip
  \gradetable[h]
\end{center}

\vspace{0.5cm}

\begin{questions}

% --- EXERCICE 1 : STRUCTURES ET DICTIONNAIRES (9,5 points) ---
\question \textbf{Maîtrise des types construits}

\begin{parts}
    \part[2] Pour chaque besoin, indiquez la structure (\textbf{Liste} ou \textbf{Dictionnaire}) la plus adaptée. Justifiez (0,5 pt par réponse). 
    \begin{itemize}
        \item Gérer le carnet d'adresses d'un téléphone (nom de contact $\rightarrow$ numéro) : \dotfill\\
        Justification : \dotfill 
        \item Enregistrer le classement des 10 meilleurs scores d'un jeu (du 1er au 10ème) : \dotfill\\
        Justification : \dotfill
        \item Associer des codes d'erreurs (404, 500) à leur message de description : \dotfill\\
        Justification : \dotfill
        \item Historique des pages web visitées dans l'ordre chronologique : \dotfill\\
        Justification : \dotfill
    \end{itemize}

    \part[1] \textbf{Création :} Écrivez une instruction Python pour créer un dictionnaire nommé \texttt{vaisseau} possédant deux clés : \texttt{"modele"} (valeur "X-Wing") et \texttt{"energie"} (valeur 80). \\
    .\dotfill

    \part[1] \textbf{Accès :} Soit \texttt{prix = \{"Pomme": 1.5, "Banane": 2.1\}}. 
    Écrivez l'instruction pour récupérer le prix de la "Banane" dans une variable \texttt{p}. \\
    .\dotfill

    \part[1] \textbf{Modification :} Écrivez l'instruction pour modifier le prix de la "Pomme" à 1.8 dans le dictionnaire précédent.\\
    .\dotfill

    \part[1] \textbf{Recherche :} Écrivez une instruction permettant de tester si la clé \texttt{"Cerise"} est présente dans \texttt{prix} (doit renvoyer un booléen).\\
    .\dotfill

    \part[\half] \textbf{Analyse de boucle :} Soit le code suivant : 
        \begin{lstlisting}[language=Python]
for v in prix.values():
    print(v)\end{lstlisting}
        Quelle catégorie d'informations (les noms des fruits, les prix ou les deux) sera affichée ?\dotfill
    
    \part[1] \textbf{Analyse de boucle :} On souhaite afficher chaque fruit suivi de son prix. 
    Complétez la ligne manquante :
    \begin{lstlisting}[language=Python]
for fruit, valeur in ................................... :
    print(fruit, "coûte", valeur)\end{lstlisting}

    \part[\half] \textbf{Analyse de boucle :} Si on utilise la boucle \texttt{for k in prix.keys():}, que contiendra la variable \texttt{k} au premier tour de boucle ? \\
    .\dotfill

   \part[2] \textbf{Parcours :} Écrivez un script (boucle \texttt{for}) qui parcourt le dictionnaire \texttt{prix} et affiche uniquement le nom des fruits dont le prix est $> 2.0$.\\
    .\dotfill\\
    .\dotfill\\
    .\dotfill
\end{parts}

\vspace{0.4cm}

% --- EXERCICE 2 : STRUCTURES COMPOSITES (3 points) ---
\question \textbf{Gestion d'une flotte de véhicules} 
\begin{lstlisting}
garage = [
    {"marque": "Renault", "km": 45000, "etat": "neuf"},
    {"marque": "Peugeot", "km": 120000, "etat": "occasion"},
    {"marque": "Tesla", "km": 5000, "etat": "neuf"}
]
\end{lstlisting}

\begin{parts}
    \part[1] Quelle instruction permet d'accéder au \textbf{kilométrage} de la Tesla ? (3ème élément)\\
    .\dotfill
    
    \part[1] \textbf{Utilisation de len() :}\\ Que renvoie l'instruction : \texttt{len(garage)} ?\dotfill\\
    Justification : \dotfill 
    
    \part[2] Complétez la fonction ci-dessous pour qu'elle renvoie le kilométrage moyen des véhicules d'un certain \textit{etat}. 
\begin{lstlisting}
def moyenne_km(liste_voitures, etat_voulu):
    total_km = 0
    compteur = 0
    for voiture in ......................... :
        if voiture["etat"] == etat_voulu:
            total_km = total_km + .........................
            compteur = compteur + 1
    return total_km / compteur
\end{lstlisting}
\end{parts}

% --- EXERCICE 3 : ALGORITHMES GLOUTONS (7,5 points) ---
\question \textbf{Le problème du cambrioleur (Sac à dos)}

Un cambrioleur possède un sac pouvant transporter au maximum \textbf{10 kg}. 
\begin{center}
\begin{tabular}{|l|c|c|c|c|}
\hline
Objet & Console & Tablette & Statuette & Caméra \\ \hline
Poids (kg) & 5 & 2 & 7 & 3 \\ \hline
Valeur (\euro) & 300 & 400 & 700 & 450 \\ \hline
\end{tabular}
\end{center}

\begin{parts}
    \part[1] \textbf{Analyse d'erreur :} Un élève tape \texttt{print(objet[0])} pour voir le nom "Console". Pourquoi est-ce correct si \texttt{objet} est une liste, mais poserait problème si c'était un dictionnaire ? \\
    .\dotfill

    \part[1] \textbf{Stratégie 1 :} En choisissant la "plus grande valeur d'abord" (\textit{Capacité : 10 kg}), quels objets sont pris et quelle est la valeur totale ?\\
    .\dotfill

    \part[1] \textbf{Stratégie 2 :} Calculez les ratios (Valeur/Poids) de la Tablette et de la Caméra. Laquelle est prioritaire selon cette stratégie ? \\
    .\dotfill

    \part[1] Trouvez une combinaison d'objets (respectant les \textbf{10 kg}) qui rapporte plus que la Stratégie 1. Justifiez.\\
    .\dotfill

    \part[2] \textbf{Synthèse :} Complétez avec : \textit{rapide, local, optimal, globale}.  \\
    L'algorithme glouton est une méthode ......................... car il ne revient jamais sur ses pas. À chaque étape, il fait le meilleur choix ......................... . Cependant, ce choix ne garantit pas d'obtenir une solution ......................... ......................... .
\end{parts}

\end{questions}

\begin{center}
    \rule{0.5\textwidth}{0.4pt}\\
    \footnotesize Fin de l'énoncé de rattrapage.
\end{center}

\end{document}